深度学习 Part 的前面单元解释模型怎样计算和训练，却没有回答这些机制在具体任务中利用了什么结构。这个单元不承诺给出``深度网络为什么有效''的统一答案，而是比较几种可检验的数学视角，并说明每种视角能解释什么、遗漏什么。

本单元回答四个问题，它们按``从理论到实践、从理解到解释''的顺序展开：

\begin{enumerate}
\tightlist
\item
  神经网络的函数空间视角：NTK 与 mean-field 来自不同的无限宽缩放，用于分析特定训练极限，而不是把所有有限网络化成同一个对象。
\item
  特征学习与核方法的差异：固定特征和随训练改变的特征分别获得什么能力，以及怎样用实验判断网络是否真的离开了近核制度。
\item
  高维数据的低维结构检验：随机投影、局部维数和流形假设分别检验什么；二维可视化为什么不能充当低维证据。
\item
  模型解释的数学形式化：SHAP、LIME 与积分梯度回答的是归因问题；归因稳定性、因果解释和机制理解为何不能互相替代。
\end{enumerate}

四篇共用一条主线：先明确分析对象和极限制度，再定义要测量的结构，最后寻找能推翻解释的对照实验。理论类比和漂亮可视化只能产生假设，不能代替证据。

阅读本单元前，建议先熟悉 Part 13 的机器学习基础、Part 12 的流形与信息几何、以及 Part 14 前面单元的深度学习技术。这些前置概念会在本单元中作为``理解深度学习''的具体例子出现。

\subsection{怎么读这一章}

四篇分别从函数空间、核极限、低维结构和解释形式化切入。不要把可视化图案直接当机制解释；始终说明解释针对哪个输入、哪个输出、哪种扰动，以及是否经过反事实或稳定性检验。

\subsection{本章篇目}

以下篇目按概念依赖排列；若从具体问题进入，也可以先打开最接近当前困惑的一篇，再沿篇内链接回补。

\begin{itemize}
\tightlist
\item \hyperref[sec:14-Deep-Learning/07-Interpretability-Mathematics/01]{``神经网络的函数空间视角''}；
\item \hyperref[sec:14-Deep-Learning/07-Interpretability-Mathematics/02]{``特征学习与核方法的差异''}；
\item \hyperref[sec:14-Deep-Learning/07-Interpretability-Mathematics/03]{``高维数据的低维结构检验''}；
\item \hyperref[sec:14-Deep-Learning/07-Interpretability-Mathematics/04]{``模型解释的数学形式化''}。
\end{itemize}
